\maketitle

\begin{abstract}
\todo[inline]{Abstract --- written last, once the application numbers
are verified. Will state: (a) a delay-embedding local-Gaussian
conditional-kernel predictor instantiated on Ontario IESO demand;
(b) the validated out-of-sample backtest; (c) the directly-verified
infeasibility of a retrospective IESO head-to-head and the registered
forward experiment built in response; (d) the structural caveat that
\(\Pij\) departs from \(\kQ\) in both location and shape, with the
recoverable loss of each quantified and shown to fall below the
pre-registered actionability bar on the full population. The method's
theoretical positioning relative to the Resolvent\_Framework programme
is summarised here and treated in depth in a companion paper.}
\end{abstract}

\section{Introduction}
\label{sec:intro}

This paper frames short-horizon Ontario electricity-demand forecasting
as a question about the \emph{conditional law of a reconstructed
dynamical state}: given a delay embedding of recent demand, what is
the distribution of the next state? We instantiate a local-Gaussian
conditional-kernel estimator of that law on IESO data and evaluate it
as a \emph{prediction experiment} --- with an out-of-sample backtest,
a hash-stamped registered forward experiment, and an explicit
accounting of what can and cannot be compared against IESO's own
published forecasts.

The estimator has a precise relationship to the conditional-regularity
kernel \(\kQ\) of the Resolvent\_Framework
programme~\cite{ResolventFramework}; the correspondence and its
numerical demonstration on controlled systems are the subject of a
companion method paper in preparation. Section~\ref{sec:method}
states only what is needed to interpret the Ontario results.

\emph{Positioning relative to prior work is deliberately deferred.}
\todo[inline]{Prior-work context, for the author group to supply; this
draft makes no comparative-novelty claims.}

\section{Theory and method}
\label{sec:method}

\subsection{Construction}
\label{sec:construction}

Let a scalar observable (here, de-seasonalised hourly demand;
Section~\ref{sec:pipeline}) be delay-embedded into \(\Rb^{d}\) by lags
\(x_t = (s_t, s_{t-1}, \ldots, s_{t-(d-1)})\). For a query state \(x\)
we form a local neighbourhood in the embedding space and fit, by
weighted least squares, a local linear one-step map
\begin{equation}
\label{eq:drift}
x_{\text{next}} \;\approx\; x \cdot \Cj \,,
\qquad\text{(row convention)}
\end{equation}
with \(\Cj \in \Rb^{d \times d}\) per anchor \(j\). Residuals
\(Y - X \Cj\) over the local neighbourhood define a diffusion
\begin{equation}
\label{eq:diff}
\Sj \;=\; \text{plain $\mu$-centred covariance of }\bigl(Y - X \Cj\bigr) \,,
\end{equation}
\emph{without} a residual kernel. The pair defines a state-dependent
Gaussian Markov kernel on the embedding space
\begin{equation}
\label{eq:kernel}
\Pij(x, \cdot) \;=\; \Nrm\bigl(x \Cj,\; \Sj\bigr).
\end{equation}

\paragraph{Locality of \(\Cj\).} The Goldenshluger--Lepski selector in
the supporting library is degenerate for the normalised Gaussian
kernel: its score is monotone in \(\theta\) with no interior optimum.
True leave-one-out CV on one-step squared error is also
non-discriminating --- under bootstrap-SE testing across anchors the
LOO objective is flat-within-noise above a small \(\theta\), with no
statistically separated interior minimum. This extends to multi-step
LOO at every horizon \(h \in \{1, 2, 4, 8, 16, 24\}\) and reproduces
on three controlled systems (linear VAR(1), logistic map,
regime-switching AR), so it is a property of the \emph{criterion}
(prediction-error CV does not identify a localisation bandwidth here),
not of the Ontario data. Consequently, \(\Cj\) is effectively a
single global linear operator at the \(d = 2\) z-score-lag embedding.
The
plain (not residual-kernel-weighted) covariance is used for \(\Sj\)
because the weighted form collapses toward \(\approx w^{2} \cdot Q\).

With \(\Cj\) effectively global and linear, the non-Dirac content of
the predictive law lives entirely in \(\Sj\) and its departure from
Gaussianity (Section~\ref{sec:caveat}). The available lever is the
\emph{state representation} (richer embedding, explicit
phase/derivative coordinate), not the bandwidth rule. Estimator
internals: Appendix~\ref{app:estimator}.

\paragraph{Relationship to the Resolvent\_Framework kernel \(\kQ\).}
The Resolvent\_Framework programme~[RF, \S``Conditional regularity and
dynamics''] defines, by Rokhlin disintegration of two observations
\(Q, F\), a Markov kernel
\begin{equation}
\label{eq:kappaQ}
\kQ(q, \cdot) \;=\; P(F \in \cdot \mid Q = q),
\end{equation}
``forced by the measure, not chosen,'' in a register the programme
labels ``Disclosure, not construction. The structures are \emph{in}
the measure.'' (\cite{ResolventFramework} lines 69--70 and 63).
Indexing over time yields a Markov semigroup \(\{\Pi_t\}\)
(Chapman--Kolmogorov \emph{derived}, not assumed) and Koopman--Perron
duality
\begin{equation}
\label{eq:KP}
\int K_t\, g\, \dd\mu \;=\; \int g\, \dd(P_t^{*}\mu);
\end{equation}
\(K_t\) collapses to the classical Koopman operator when the kernels
are Dirac. Three layers are kept notationally distinct: \(\kQ\)
(disintegration), \(\{\Pi_t\}\) (semigroup), \(K_t\) (operator).

The object Section~\ref{sec:construction} constructs is a
finite-sample, parametric, single-step, Gaussian-form estimator of
\(\kQ\) at the embedding layer, from which the operator-layer objects
\(K_{\Delta} / P_{\Delta}^{*}\) are recoverable in the non-Dirac
regime (\(\Sj \neq 0\)). The qualifier ``locally'' is dropped: per
Section~\ref{sec:construction} no statistically meaningful locality
scale is recoverable at this embedding; the drift reduces to a single
global linear operator. The \(\kQ\) correspondence holds at the
embedding layer; \emph{locality} of that estimator is not empirically
operative on Ontario. The delay embedding corresponds to the
programme's reconstruction object; \(\Cj\) to a linearisation of the
minimal-sufficient factor (intended local, effectively global here);
\(\Sj\) to the second moment of \(\kQ\)'s Gaussian projection.
\(\Sj \to 0\) reproduces the classical-Koopman collapse. Three
qualifiers bound this correspondence:

\begin{enumerate}
  \item \textbf{Register.} The programme \emph{discloses} \(\kQ\) from
  the measure; this work \emph{constructs} a Gaussian proxy and fits
  it from finite data --- ``estimator of,'' not ``instance of.''

  \item \textbf{Imposed Gaussianity.} \(\Pij = \Nrm(x \Cj, \Sj)\)
  coincides with \(\kQ\) only where the local conditional law is
  Gaussian. Section~\ref{sec:caveat} quantifies the gap as real but
  below the pre-registered actionability bar.

  \item \textbf{Single step.} \(\Pij\) is the single-step slice
  \(\Pidelta\); the semigroup \(\{\Pi_t\}\) is neither composed nor
  verified here. (A multi-step predictive-variance propagation via a
  shift-aware state Jacobian is built and gate-validated as the
  instrument for Section~\ref{sec:caveat}'s decomposition; it is not
  the semigroup.)
\end{enumerate}

Numerical demonstration on controlled systems is the subject of the
companion method paper.

\section{Structural caveat: a mis-centred conditional mean and a
Gaussian-proxy diffusion}
\label{sec:caveat}

The two estimator factors \(\Cj\) and \(\Sj\) each carry a quantified
limitation. A pre-registered dev-set decomposition (below), on the
full population with predictive variance correctly propagated, finds
\textbf{neither effect clears the materiality bar}.

\paragraph{\(\Sj\) is rank-1 by construction.} For a single-variable
delay embedding, coordinates \(2, \ldots, d\) of the one-step image
are deterministic shifts of the input coordinates
(\(Y_{[:,\,1:]} = X_{[:,\,:-1]}\) exactly). The only stochastic
content of the one-step map is the scalar coordinate-0 innovation;
therefore \(\Sj\) has exactly rank one, independent of data --- the
multivariate diffusion is structurally a single scalar innovation
variance embedded in \(\Rb^{d \times d}\). Any ``multi-mode'' reading
of \(\Sj\) is incoherent for this embedding and is not reported.

\paragraph{That scalar innovation is non-Gaussian.} The
non-Gaussianity question is therefore \emph{univariate}. Through the
production estimator, validated against synthetic ground truth
(Section~\ref{sec:synth}), the one-step innovation is heavy-tailed;
the propagation-aware strength of that departure
(Section~\ref{sec:caveat-propagation}) is what the diagnostic of
Section~\ref{sec:caveat-decomposition} acts on.

\paragraph{The conditional mean is phase-mis-centred.} The local
linear map cannot distinguish a state rising through a demand level
from one falling through it (the embedding lacks a phase coordinate),
so it centres between the two regimes. The signed \(z\)-space error
correlates with \(\dd z/\dd t\) at \(\mathrm{corr} = -0.246\)
(bootstrap 95\% CI \([-0.283,\, -0.207]\), \(n = 3167\)), monotone
across rising/falling phase. \(\Pij\)'s location, not only its shape,
departs from \(\kQ\).

\paragraph{Decomposition.}
\label{sec:caveat-decomposition}
A pre-registered dev-set diagnostic (2021--22 inspection benchmark;
inspection-only, not claim-grade) bounds the recoverable loss of each
factor via oracle counterfactuals on multi-step forecast log-loss
over all delivery days, with predictive variance propagated across
the embedding-dimension change (below). The phase-bias oracle
recovers \(\approx 4.8\%\) of dev log-loss (95\% CI
\([0.7\%,\, 8.6\%]\)); the Student-\(t\) (\(\nu \approx 6.5\))
innovation oracle recovers \(\approx 2.3\%\) (95\% CI
\([-0.0\%,\, 4.8\%]\)). \textbf{Neither lower bound clears the
pre-registered 1\% materiality threshold, and the heavy-tail interval
includes zero.} The verdict is the \emph{catch-all} ``neither
dominates'' --- not because the effects are real and tied (which was
the single-dimension subpopulation case, where both cleared the bar
at \(\approx 6\%\) / \(\approx 9\%\)), but because on the complete,
correctly-propagated data neither is established at all. Document
both, engineer neither.

\paragraph{Propagation reconciliation.}
\label{sec:caveat-propagation}
The heavy-tail signal attenuated roughly four-fold
(\(\nu \approx 5.0 \to 6.5\); implied excess kurtosis
\(\approx 6.1 \to 2.4\)) once predictive variance was propagated
across the dimension-change boundary rather than sidestepped by
excluding those days. Earlier one-step measurements
(re-baseline excess kurtosis \(\approx 24\)--\(33\), without
multi-step variance propagation) were partly a propagation artefact,
not solely a property of the innovation law. Propagation is supplied
by a fixed-maximum-dimension state augmentation carrying unused lag
coordinates as deterministic, zero-innovation components, validated
to machine precision against the closed-form linear-Gaussian
reference and bit-identical to the unaugmented recursion on
single-dimension days. Restricting to the single-dimension
subpopulation (\(\approx 58\%\) of days) yields the same qualitative
verdict.

\section{Application and results}
\label{sec:application}

\subsection{Ontario pipeline}
\label{sec:pipeline}

The instantiation is a configuration-driven pipeline; all paths and
parameters resolve from a single \texttt{config/pipeline.yaml}
relative to the project root (no working-directory dependence).
Stages: \textbf{acquire} IESO public \texttt{PUB\_Demand};
\textbf{de-seasonalise} (Section~\ref{sec:zscore}) hourly demand to a
stationary \(z\)-score; \textbf{day-type tag}
\(\{\mathrm{weekday}, \mathrm{saturday}, \mathrm{sunday}\}\) using a
configurable 07:00 day-anchor offset (transitions whose one-step
target crosses the day-type rollover seam are excluded --- without
this \(\Sj\) is catastrophically ill-conditioned, so the estimand is
effectively forced to be \emph{intra-day}); \textbf{embedding
dimension} selected per day-type by an elbow rule on the
prediction-skill-vs-dimension curve; \textbf{per-anchor fit} of the
Section~\ref{sec:method} estimator.

\subsection{The \(z\)-score transform is a framework coupling, not a
preprocessing detail}
\label{sec:zscore}

De-seasonalisation is
\begin{equation}
\label{eq:zscore}
z \;=\; \frac{D - \mu_{m,h}}{\sigma_{m,h}},
\end{equation}
with \(\mu, \sigma\) a month \(\times\) hour-of-day climatology
estimated strictly from pre-cutoff data. This single climatology
object is woven through \emph{all four} estimator layers: the embedded
state, the drift fit, the diffusion (\(\Sj\) is the covariance of
\(z\)-residuals), and the demand-space round-trip. It is therefore an
implicit seasonal model coupled into \(\kQ\), not a benign
preprocessing step.

Re-running under climatology specs of varying resolution (including
no-transform) tested whether the coupling corrupts the dynamics
(leaky implicit seasonal model --- ``Reading 2'') or is a benign
invertible coordinate change (``Reading 1''). The phase/asymmetry
structure was \textbf{invariant} across specs --- Reading 1 for the
phase axis: the conditional-mean structure is genuine demand
dynamics. The tail structure \emph{did} move with seasonal resolution
--- a partial Reading 2 confined to the heavy-tail axis, which is the
Section~\ref{sec:caveat} caveat, not a separate defect.
\todo[inline]{Development-set investigation; its role here is
methodological --- establishing the coordinate is sound before any
state-representation honing.}

\subsection{Out-of-sample backtest}
\label{sec:backtest}

The iterated one-step predictor is backtested on Ontario ground truth
for complete delivery days strictly after the model-specification
cutoff. Leakage guard: fits and \(z\)-score climatology use only
\(\le\)-cutoff data; post-cutoff actuals enter only as query history
and scoring truth. Error is reported per horizon \(h = 1, \ldots, 24\);
in a day-ahead forecast, horizon is collinear with hour-of-day, so
the error structure is \emph{diurnal} (deep night easiest; dawn/dusk
ramps hardest), not pure horizon compounding. The committed
provenanced C1 artifact gives \(\mathrm{MAE} \approx 786~\mathrm{MW}\),
\(\mathrm{MAPE} \approx 4.78\%\) over 500 out-of-sample days,
\(\approx 1.5\times\) better than seasonal persistence at every hour.
\todo[inline]{Current CLAIM-grade values from the committed
provenanced artifact; the artifact hash will be restated after any
honing re-freeze. Per-horizon table will appear in
Appendix~\ref{app:backtest}.}
Interval coverage carries the Section~\ref{sec:caveat} Gaussian-proxy
caveat: a Gaussian interval around a heavy-tailed innovation
under-covers.

\subsection{Registered forward prediction experiment}
\label{sec:registered}

A fair retrospective IESO comparison is not recoverable
(Section~\ref{sec:ieso}), so the comparison is run forward. A
hash-stamped frozen specification captures estimator identity,
resolved hyperparameters, data cutoff, code commit, and a content
fingerprint of pre-cutoff actuals, and is verified on load. The
predictor is local/lazy (per-anchor fit at prediction time from
\(\le\)-cutoff library data), so the freeze pins the function and
its inputs, and a cutoff guard blocks post-cutoff data from any fit.
Forecasts issue forward; scraped actuals settle into an append-only
ledger (IESO restates recent demand --- revisions appended, never
overwritten); scoring runs against persistence baselines with IESO
published values recorded as differently-horizoned context.

\todo[inline]{Forward-experiment results accrue with calendar time;
report once a sufficient settled window exists.}

\subsection{The retrospective IESO head-to-head is infeasible from
public archives}
\label{sec:ieso}

Full enumeration of all 133 public IESO report directories establishes
that no public archive is \emph{both} Ontario-demand-basis \emph{and}
day-ahead-issued for historical dates. Two independent blockers: the
Ontario-demand-basis product (Adequacy3 \texttt{ForecastOntDemand})
retains only the end-of-delivery-day version for days old enough to
have settled actuals (wrong horizon); the day-ahead-issued products
(DATotals, PredispTotals) publish only market totals, where ``Total
Load'' carries a definitional \(\approx +2485~\mathrm{MW}\) offset
over Ontario Demand (MarketQuantity enumeration: Total
Energy/Loss/Load/Dispatchable/10S/10N/30R; none is Ontario demand).
This is \emph{why} the registered forward experiment
(Section~\ref{sec:registered}) is the methodologically correct
comparison. (Stated for public archives as of the recorded access
date.)

\subsection{Validation against synthetic ground truth}
\label{sec:synth}

Three gates call the \emph{production} code path (not
reimplementations):

\begin{enumerate}
  \item \textbf{Recovery} --- a VAR(1) system with known \((A, Q)\) is
  recovered within tolerance (drift and diffusion).
  \item \textbf{Non-Gaussianity} --- a Gaussian-innovation VAR(1)
  reads approximately Gaussian; a Student-\(t\)-innovation VAR(1) is
  flagged heavy-tailed, validating the Section~\ref{sec:caveat}
  diagnostic against ground truth.
  \item \textbf{Multi-step composition} --- iterated \(\Pidelta\) is
  checked against the VAR(1) closed form; drift error compounds with
  horizon \emph{by construction} and is reported as an error-growth
  curve within tolerance.
\end{enumerate}

All three gates pass.
\todo[inline]{Cite the provenanced gate artifact + its numbers from
the verified run.}

\section{Limitations and reproducibility}
\label{sec:limits}

\begin{itemize}
  \item \textbf{Estimator limitations (Section~\ref{sec:caveat}).}
  \(\Pij\) departs from \(\kQ\) in location (\(\Cj\)
  phase-mis-centred) and shape (\(\Sj\) Gaussian proxy of a
  heavy-tailed innovation). The pre-registered decomposition bounds
  recoverable loss at \(\approx 4.8\%\) and \(\approx 2.3\%\) of dev
  log-loss; \textbf{neither clears the pre-registered 1\% materiality
  bar} (heavy-tail interval includes zero). The heavy-tail effect
  attenuated \(\approx 4\times\) under propagated variance, indicating
  it was partly a propagation artefact. Open: phase-aware conditional
  mean and/or non-Gaussian local model; neither licensed as the
  higher-leverage lever by current evidence.

  \item \textbf{Single-step only.} The semigroup \(\{\Pi_t\}\) is not
  constructed; multi-step is iteration of \(\Pidelta\), validated for
  error growth. (The companion method paper is the natural home for
  the semigroup treatment.)

  \item \textbf{Exogenous information.} No weather/exogenous
  covariates; the backtest error structure suggests this is a
  model-scope ceiling, not a hyperparameter issue.

  \item \textbf{[DEFER-RF] points.}
  \todo[inline]{Residual theory-correspondence questions for resolution
  with the Resolvent\_Framework authors; several are mooted by the
  rank-1 finding.}
\end{itemize}

\paragraph{Provenance.} Every claim- or method-grade result is written
through a single envelope stamping each artifact with two hashes ---
an inputs fingerprint (git commit + frozen-spec hash + library
versions + RNG seeds + declared inputs) and a body hash recomputed on
read --- refuses to write from a dirty source tree, and carries a
mandatory grade banner (CLAIM / METHOD / INSPECTION-ONLY).
Predictor-derived claim artifacts additionally bind the frozen-spec
hash. Inspection-only material is banner-marked and not cited as a
result.

\section{Conclusion}
\label{sec:conclusion}

The two ways \(\Pij\) departs from \(\kQ\) --- phase-mis-centred mean
and Gaussian-proxy diffusion --- were pre-registered for elimination
and both fall below the materiality bar on the validated full
population. The standing
shortcomings are structural (univariate state, phase-blind embedding,
no exogenous covariates), not tuning issues, and the higher-leverage
next step is a state representation change rather than an estimator
change.

\appendix

\section{Estimator detail}
\label{app:estimator}

\todo[inline]{WLS weighting, the LOO-CV \(\theta\) rule, the
GL-degeneracy note in full. The companion method paper is the primary
home for this; reproduce here only what the application reader needs.}

\section{Per-horizon backtest table}
\label{app:backtest}

\todo[inline]{From the C1 artifact.}

\section{Provenance requirements matrix}
\label{app:provenance}

Reference: \texttt{experiment/results/PROVENANCE\_REQUIREMENTS.md}.


% ---------- References ----------
\ifstandalone
\bibliographystyle{plain}
\bibliography{references}
\fi

% (Additional citations: \nocite{CompanionMethodPaper} ensures the
% companion-method-paper entry appears in the bibliography even before
% it is cited in the body; remove the \nocite once a body \cite{} is
% in place.)
\nocite{CompanionMethodPaper}
